\section{Discusión}
\label{sec:discusion}

Organizamos la discusión en cuatro hallazgos. Para cada uno decimos qué encontramos, qué implica para la decisión de
Ternium, qué no implica y qué dato o experimento cerraría la brecha. Los cuatro viven en los dos primeros peldaños de
la escalera de afirmaciones: \emph{observamos} y \emph{se asocia, ajustando por}.

\subsection{H1. La corrección depende más del programa, o del periodo, que del producto}
\emph{Qué encontramos.} Por encima del azar, el programa del Skinpass explica \ExcesoPrograma~puntos de la variación de
la corrección de bending \ExcesoProgramaIC, el producto \ExcesoProducto{} \ExcesoProductoIC{} y la dimensión
\ExcesoDimension; el predominio del programa se mantiene si entra antes que el producto y en las \NVariantes{}
alternativas de depuración y agrupación y agrupación. Agrupar por día da casi lo mismo que por programa (\VCuatroExcesoDia{} puntos).

\emph{Qué implica.} Una receta por producto podría absorber, como máximo, la parte sistemática de la corrección que
viaja con el producto: \ExcesoProducto{} puntos de su suma de cuadrados por encima del azar (\ExcesoProdDim{} con la
dimensión). La parte más grande cambia de un programa a otro o de un día a otro, y ninguna receta por producto la
absorbería. Para decidir, lo que más pesa es qué cambia entre programas: el rodillo de trabajo, el operador, la mezcla
de pedidos del periodo o el estado de los medidores.

\emph{Qué no implica.} No implica que el programa cause la corrección: con estos datos, el programa y el día son casi
indistinguibles, por lo que escribimos ``programa o periodo''. Tampoco implica que los productos respondan distinto a la
receta. Que un producto necesite en promedio una corrección distinta es una diferencia de \emph{nivel}; que necesite
una receta distinta exige una diferencia de \emph{respuesta}, es decir, que el mismo cambio de receta tenga efectos
distintos en productos distintos, y eso no se observa en datos donde la receta casi no varía dentro de producto. Por
último, no implica que reducir la corrección mejore la calidad: el Skinpass no registra planeza de salida.

\emph{Qué lo cerraría.} El registro, por rollo, de la campaña de rodillos de trabajo (hoy solo el identificador del
rodillo), de los cambios de receta con su fecha de vigencia y de los ajustes manuales del operador con su motivo
permitiría separar el programa del periodo. El experimento de la sección~\ref{sec:experimento} bloquea ese efecto en
lugar de estimarlo.

\subsection{H2. La corona del MC4 se asocia con la corrección, pero poco y de forma frágil}
\emph{Qué encontramos.} Dentro de producto, dimensión y mes, un punto porcentual más de corona relativa se asocia con
una corrección \PendCorona~kN más negativa \PendCoronaIC, con el signo que predice el argumento mecánico. El control
negativo no responde, la dirección sobrevive a \NVariantes{} alternativas y el preset no reacciona a la corona. Pero el
efecto es pequeño: del percentil 10 al 90 de la corona la corrección cambia \VDosCambio~kN, frente a una desviación
estándar de \YDEAsoc~kN. Además es frágil: un confusor no medido con 1/\CHUnoEntre{} de la fuerza del producto bastaría
para anularlo.

\emph{Qué implica.} La forma con que el rollo sale del MC4 lleva poca información útil para fijar el preset del
Skinpass. Ajustar el preset por la corona movería la corrección unos pocos kN, una fracción pequeña del trabajo que
hoy hace el lazo (\ExpAbsY~kN de $|\Y|$ promedio).

\emph{Qué no implica.} No es un efecto causal, y la fragilidad lo impide aún más que de costumbre: la planeza de
entrada y el estado de los rodillos del MC4 influyen tanto en la corona como en lo que el lazo necesita corregir. Tampoco
implica que la corona no importe para la planeza final: el lazo compensa la forma de entrada, y lo que medimos es cuánto
la compensó, no si la compensación fue suficiente.

\emph{Qué lo cerraría.} Una medición de planeza a la entrada del Skinpass, que hoy no existe, convertiría a la principal
variable no medida del diagrama en una variable observada.

\subsection{H3. No se detectan diferencias de respuesta entre productos}
\emph{Qué encontramos.} De los \NProdHet{} productos con al menos 300 rollos, \NProdNeg{} tienen una pendiente
corona~$\rightarrow$~$\Y$ negativa como la global, pero ninguno se distingue de cero tras controlar la tasa de
descubrimientos falsos y los intervalos se traslapan ampliamente.

\emph{Qué implica.} Con estos datos no hay evidencia de que algún producto reaccione a la forma de entrada de manera
distinta a los demás. Es la única forma de interacción que podemos estimar, y no aparece.

\emph{Qué no implica.} No demuestra que los productos respondan igual: con intervalos de este ancho, la ausencia de
evidencia no es evidencia de equivalencia, que exigiría fijar un margen de equivalencia de antemano~\cite{wellek2010}.
Los productos de los extremos de la Fig.~\ref{fig:productos} son los peores candidatos para una receta especial: con
\NProdHet{} comparaciones, el más extremo sobrestima su efecto por la maldición del optimizador~\cite{smith2006}.

\emph{Qué lo cerraría.} Variación deliberada del preset dentro de producto, que solo puede venir de un experimento.

\subsection{H4. La información previa al Skinpass no anticipa la corrección mejor que una receta por producto}
\emph{Qué encontramos.} El modelo de gradient boosting baja el error absoluto medio \PredMAEDifGrupo~kN respecto a la
media por producto y dimensión en la validación agrupada, pero es peor en la temporal, y los dos explican menos del
10\,\% de la variación fuera de muestra. Lo que el modelo usa es, sobre todo, el producto. Recentrar el preset con la
mediana histórica bajaría $|\Y|$ unos \ExpDeltaA~kN; casi todo eso (\ExpDeltaGlobal~kN) se logra con una sola mediana
global, y diferenciar por producto agregaría \ExpDeltaDif~kN.

\emph{Qué implica.} Optimizar el preset con un modelo más flexible sobre la información previa al Skinpass agregaría muy
poco a una receta por producto, y una receta por producto agregaría muy poco a recentrar la receta única. Lo que sí
sugieren los datos es que la receta vigente está desplazada (la corrección mediana es de \YMediana~kN), y esa es una
hipótesis operativa que vale la pena probar.

\emph{Qué no implica.} No implica que la corrección sea impredecible con otros datos; el modelo no conoce la planeza de
entrada ni el estado de los rodillos. Tampoco implica que un preset recentrado mejore la calidad: el lazo corrige de
todos modos, y la ganancia, si la hay, estaría en el tramo de cabeza del rollo, antes de que el lazo converja. Los datos
resumidos por rollo no separan ese tramo.

\emph{Qué lo cerraría.} El experimento de la sección~\ref{sec:experimento}, con planeza de salida medida durante su
ejecución.

\subsection{Relación con la literatura}
Los resultados coinciden con la advertencia de Box~\cite{box1966}: las variables latentes del proceso, aquí resumidas
en el programa o el periodo, pesan más que las variables observadas que solemos poner en un modelo, y la asociación de
la corona sería la conclusión equivocada de un análisis que no la sometiera a sensibilidad. También coinciden con la
física del temple~\cite{lenard2014} en que la forma de entrada importa (el signo de la asociación es el esperado),
pero matizan su alcance práctico: en una planta con un lazo de planeza, esa forma deja una huella pequeña en la
corrección. Y chocan con la expectativa, común en la literatura de aprendizaje automático en manufactura~\cite{wuest2016},
de que un modelo más flexible extraerá de los datos de proceso lo que una regla simple no ve: con validación por
programa y por tiempo, la ventaja casi desaparece, como advierten Roberts et al.~\cite{roberts2017}.

\subsection{Limitaciones, escritas como hallazgos}
\begin{itemize}
\item \emph{Sin resultado de forma.} El Skinpass no registra planeza. Sus defectos de inspector afectan al
      \PctDefInspector\,\% de los rollos y el \PctTensionEnrollado\,\% de esos registros son de tensión de enrollado, no
      de forma; por eso no pudimos usar la calidad de forma como resultado.
\item \emph{Censura de fin de año.} Solo \PctDicEnSKP\,\% de los rollos de diciembre llega al Skinpass de 2025, y los de
      diciembre traen más corona (\CoronaRelMedianaDic\,\% contra \CoronaRelMediana\,\%). Sin el corte, la pendiente se
      acerca a cero y su intervalo lo incluye.
\item \emph{La receta no es tan única.} El preset es único en el \PctPresetUnicoTodos\,\% de los grupos programa
      $\times$ producto $\times$ dimensión, pero solo en el \PctPresetUnicoDos\,\% de los que tienen dos o más rollos;
      \PctPresetCero\,\% de los rollos tiene preset cero y \NPresetAlto{} rollos tienen un preset mayor que el máximo
      bending real observado. Las unidades del preset son un supuesto (Anexo~\ref{sec:anexo-supuestos}).
\item \emph{Programa y periodo, indistinguibles.} El día explica \VCuatroExcesoDia{} puntos contra
      \VCuatroExcesoProg{} del programa.
\item \emph{Confusión fuerte por pedido.} El producto explica \CHRDRef\,\% de la varianza de la corona dentro de
      dimensión y mes; cualquier confusor parecido, aunque mucho más débil, cambiaría la conclusión sobre la corona.
\item \emph{Un año, una planta, un resumen por rollo.} No sabemos qué resume cada valor (promedio, punta u otro) ni
      separamos el enhebrado del tramo estable.
\end{itemize}
